% =========================================================================
% Capítulo 3: Cinética de sistemas de partículas
% Tema 01: Sistemas de partículas
% Problema 07: Barra giratoria en guía semicircular con resortes
% =========================================================================
\subsubsection{Barra giratoria en guía semicircular con resortes}

Un alambre con forma de semicircunferencia de radio $R$ se encuentra fijo en el plano vertical. Sobre este alambre desliza el extremo $P$ de una barra rígida ideal de masa nula, la cual puede girar libremente alrededor de un eje fijo en el centro geométrico $O$. La barra sostiene dos partículas: una de masa $2m$ en el extremo $P$ y otra de masa $m$ en el extremo opuesto $Q$. La distancia entre $O$ y $P$ es $R$, mientras que la distancia entre $O$ y $Q$ es $d$. Además, la partícula $P$ está conectada a dos resortes idénticos de constante elástica $k$ y longitud natural $l_0$, los cuales siguen la curvatura del alambre y están anclados en los extremos de su diámetro horizontal. Considerando que la gravedad actúa hacia abajo y definiendo $\theta$ como el ángulo que forma el segmento $OP$ con la vertical ascendente, realice lo siguiente:

\begin{enumerate}
    \item Defina los marcos de referencia y sistemas de coordenadas adecuados.
    \item Analice la cinemática del sistema para expresar la velocidad de las partículas $P$ y $Q$ en función del ángulo $\theta$.
    \item Analice la cinética del sistema para determinar la ecuación diferencial del movimiento utilizando métodos de energía o fuerzas.
\end{enumerate}

\begin{figure}[!ht]
    \centering
    \begin{tikzpicture}[>=Latex, scale=1.5]
        % --- Parámetros ---
        \def\R{2.2}       % Radio del semicírculo (R)
        \def\distd{1.4}   % Distancia d (OQ)
        \def\ang{25}      % Ángulo theta (desviación de la vertical)
        
        % Coordenadas Clave
        \coordinate (O) at (0,0);
        % P en coordenadas polares (90 - theta)
        \coordinate (P) at ({90-\ang}:\R);
        % Q en coordenadas polares (90 - theta + 180) -> (270 - theta)
        \coordinate (Q) at ({270-\ang}:\distd);
        % Extremos del diámetro
        \coordinate (BaseL) at (-\R, 0);
        \coordinate (BaseR) at (\R, 0);

        % --- 1. Soportes y Línea Base ---
        \draw[dashed, gray] (BaseL) -- (BaseR);
        
        % Soportes en los extremos
        \foreach \x in {-\R, \R} {
            \fill[pattern=north east lines] (\x-0.2, -0.1) rectangle (\x+0.2, 0);
            \draw[thick] (\x-0.2, 0) -- (\x+0.2, 0);
        }

        % --- 2. Resortes en la Guía (Arcos) ---
        % Resorte Izquierdo (desde 180 hasta P)
        \draw[decorate, decoration={coil, aspect=0.4, segment length=3pt, amplitude=2.5pt}] 
            (BaseL) arc (180:{90-\ang}:\R);
        % Resorte Derecho (desde 0 hasta P)
        \draw[decorate, decoration={coil, aspect=0.4, segment length=3pt, amplitude=2.5pt}] 
            (BaseR) arc (0:{90-\ang}:\R);
            
        % Línea guía del alambre (borde fino para definir la forma)
        \draw[thin, gray] (\R,0) arc (0:180:\R);

        % --- 3. Barra y Masas ---
        \draw[line width=2pt, gray!50] (Q) -- (P);
        \draw[thick] (Q) -- (P);

        % Masa P (2m)
        \shadedraw[ball color=gray!30] (P) circle (0.22) node[above right, xshift=5pt] {$P$};
        \node[font=\footnotesize] at (P) {$2m$};

        % Masa Q (m)
        \shadedraw[ball color=gray!30] (Q) circle (0.18) node[below left, xshift =-5pt] {$Q$};
        \node[font=\footnotesize] at (Q) {$m$};

        % --- 4. Cotas y Referencias ---
        % Centro O
        \filldraw (O) circle (1.5pt) node[above right, font=\small, xshift=3pt] {$O$};
        
        % Eje Vertical de Referencia
        \draw[dashed] (0, 0.5) -- (0, \R+0.3);
        
        % Ángulo theta
        \draw[->] (0, 0.8) arc (90:{90-\ang}:0.8);
        \node at (0.25, 1.0) {$\theta$};
        
        % Etiqueta R (Radio)
        \draw[thin, ->] (O) -- (135:\R) node[midway, fill=white, inner sep=1pt] {$R$};
        
        % Etiqueta d
        %\draw[thin, |-|] ($(O)!0.2!(90:\distd)$) -- ($(Q)!0.2!(90:\distd)$); 
        % Ajuste manual de cota d para que sea paralela a la barra
        \draw[thin, |<->|] ($(O)!0.15!({270-\ang+90}:1)$) -- ($(Q)!0.15!({270-\ang+90}:1)$) 
            node[midway, below right] {$d$};

        % Gravedad
        \draw[->, thick] (1.3, 1.0) -- (1.3, 0.1) node[midway, right] {$\vec{g}$};

    \end{tikzpicture}
    \caption{Sistema de barra rotatoria con masas en los extremos y resortes sobre guía semicircular.}
    \label{fig:barra_semicirculo_resortes}
\end{figure}